% !TEX root = ../main.tex
\corrigesChapitre{24}{Lunette astronomique}
\corrigesSouples

\begin{blocLivre}{cor-24-1}
\begin{corrige}{1}{Le vocabulaire de la lunette}
\demarche{Objectif : grande distance focale ; oculaire : petite distance focale ; afocale : $F'_1=F_2$.}
\begin{questions}
\item L'objectif est la lentille de \res{\qty{900}{mm}}, l'oculaire celle de \res{\qty{20}{mm}}.
\item Le foyer image de l'objectif est confondu avec le foyer objet de l'oculaire : l'image d'un objet à l'infini est
\res{\text{à l'infini}}, et l'œil l'observe sans accommoder.
\item $f'_1+f'_2=\res{\qty{920}{mm}}$.
\item $G=\dfrac{900}{20}=\res{45}$.
\end{questions}
\end{corrige}
\end{blocLivre}
\begin{blocLivre}{cor-24-2}
\begin{corrige}{2}{L'image de la Lune}
\demarche{Image dans le plan focal image ; diamètre $f'\tan\theta$.}
\begin{questions}
\item Dans le \res{\text{plan focal image}} de la lentille ; elle est \res{\text{réelle et renversée}}.
\item $1000\times\tan0{,}52^\circ=\res{\qty{9,1}{mm}}$.
\item \res{\qty{2,7}{mm}} avec le téléobjectif ; \res{\qty{18,2}{mm}} avec la lunette.
\item L'image est proportionnelle à la distance focale : une grande focale donne une image \res{\text{plus grande}}, donc plus
détaillée.
\end{questions}
\end{corrige}
\end{blocLivre}
\begin{blocLivre}{cor-24-3}
\begin{corrige}{3}{Changer d'oculaire}
\demarche{$G=f'_1/f'_2$ avec $f'_1=\qty{1000}{mm}$.}
\begin{questions}
\item $\res{G=\dfrac{f'_1}{f'_2}}$.
\item \res{25} ; \res{50} ; \res{111}.
\item La Lune entière : l'oculaire de \res{\qty{40}{mm}} (faible grossissement, grand champ). Un cratère : celui de
\res{\qty{9}{mm}}.
\item \res{\text{Oui}} : la lunette doit rester afocale, le foyer objet du nouvel oculaire doit coïncider avec $F'_1$.
\end{questions}
\end{corrige}
\end{blocLivre}
\begin{blocLivre}{cor-24-4}
\begin{corrige}{4}{Construire le faisceau émergent}
\demarche{Le rayon passant par $O_1$ n'est pas dévié ; les rayons émergents sont parallèles à $B_1O_2$.}
\begin{questions}
\item Le rayon incident qui passe par $O_1$ traverse la lentille sans déviation et coupe le plan focal image en
\res{B_1}, où convergent tous les rayons du faisceau.
\item Les trois rayons se croisent en $B_1$, puis divergent jusqu'à l'oculaire.
\item On trace le rayon $B_1O_2$ (en pointillés) : $B_1$ étant dans le plan focal objet de l'oculaire, les rayons émergents
lui sont \res{\text{parallèles}}.
\item $\tan\theta=0{,}12$ et $\tan\theta'=0{,}30$ : $\dfrac{\theta'}{\theta}\approx\res{2{,}5}=\dfrac{5}{2}=\dfrac{f'_1}{f'_2}$.
\end{questions}
\schema[corrige]{24-ex04}
\end{corrige}
\end{blocLivre}
\begin{blocLivre}{cor-24-5}
\begin{corrige}{5}{Vergences et distances focales}
\demarche{$V=1/f'$, avec $f'$ en mètres.}
\begin{questions}
\item $f'_1=\dfrac{1}{1{,}25}=\res{\qty{0,80}{m}}$ et $f'_2=\dfrac{1}{40}=\res{\qty{2,5}{cm}}$.
\item $G=\dfrac{0{,}80}{0{,}025}=\res{32}$.
\item $80+2{,}5=\res{\qty{82,5}{cm}}$.
\item $G=\dfrac{f'_1}{f'_2}=\res{\dfrac{V_2}{V_1}}=\dfrac{40}{1{,}25}=32$.
\end{questions}
\end{corrige}
\end{blocLivre}
\begin{blocLivre}{cor-24-6}
\begin{corrige}{6}{La Lune dans la lunette}
\demarche{$G=f'_1/f'_2$ ; angle exact : $\tan\theta'=G\tan\theta$.}
\begin{questions}
\item $G=\dfrac{700}{10}=\res{70}$.
\item Approximation : $70\times0{,}52^\circ=\res{36{,}4^\circ}$ ; valeur exacte : $\theta'=\arctan(70\tan0{,}52^\circ)=\res{32{,}4^\circ}$. Pour un
angle aussi grand, l'approximation des petits angles \res{\text{n'est plus valable}}.
\item $2f'_1\tan(0{,}26^\circ)=\res{\qty{6,35}{mm}}$.
\item $32{,}4^\circ<50^\circ$ : la Lune entière est \res{\text{visible}} dans le champ.
\end{questions}
\end{corrige}
\end{blocLivre}
\begin{blocLivre}{cor-24-7}
\begin{corrige}{7}{Voir les cratères lunaires}
\demarche{Taille d'un détail $=$ angle $\times$ distance ; la lunette divise par $G$ l'angle minimal à l'œil nu.}
\begin{questions}
\item $\num{3e-4}\times\num{3.84e8}=\qty{1,15e5}{m}$, soit \res{\qty{115}{km}}.
\item L'image doit être vue sous \qty{3e-4}{rad}, soit un angle réel de $\dfrac{\num{3e-4}}{70}=\qty{4,3e-6}{rad}$ : cratères de
\res{\qty{1,6}{km}}.
\item $\theta_\mathrm{min}=1{,}22\times\dfrac{\num{550e-9}}{0{,}070}=\qty{9,6e-6}{rad}$, soit \res{\qty{3,7}{km}} sur la Lune.
\item La diffraction par l'objectif fixe la finesse des détails : au-delà, augmenter $G$ agrandit une image
\res{\text{floue}} sans faire apparaître de nouveaux détails.
\end{questions}
\end{corrige}
\end{blocLivre}
\begin{blocLivre}{cor-24-8}
\begin{corrige}{8}{Programmer le grossissement}
\demarche{Image intermédiaire de taille $f'_1\tan\theta$, vue de l'oculaire sous $\theta'$ tel que $\tan\theta'=\dfrac{f'_1\tan\theta}{f'_2}$.}
\begin{questions}
\item Les fonctions trigonométriques de Python travaillent \res{\text{en radians}}. La ligne calcule l'angle $\theta'$ sous lequel
l'image est vue dans la lunette.
\item $A_1B_1=f'_1\tan\theta$ et $\tan\theta'=\dfrac{A_1B_1}{f'_2}=G\tan\theta$ : $\res{\theta'=\arctan(G\tan\theta)}$.
\item On ajoute \code{print(round(degrees(G * theta), 1))} dans la boucle. L'écart reste inférieur à 5~\% jusqu'à
$\res{G=43}$ environ.
\item La ligne \code{theta = radians(0.52)} devient \code{theta = radians(45 / 3600)}.
\end{questions}
\end{corrige}
\end{blocLivre}
\begin{blocLivre}{cor-24-9}
\begin{corrige}{9}{Retrouver les distances focales}
\demarche{$f'_1+f'_2=\qty{1,05}{m}$ et $f'_1=20\,f'_2$.}
\begin{questions}
\item $\res{f'_1+f'_2=\qty{1,05}{m}}$ et $\res{f'_1/f'_2=20}$.
\item $21\,f'_2=1{,}05$ : $f'_2=\res{\qty{5,0}{cm}}$ et $f'_1=\res{\qty{1,00}{m}}$.
\item $V_1=\res{\num{1,0}~\delta}$ et $V_2=\res{\num{20}~\delta}$.
\item $G=\dfrac{100}{2{,}5}=\res{40}$ ; il faut régler la distance entre les lentilles à \res{\qty{102,5}{cm}}.
\end{questions}
\end{corrige}
\end{blocLivre}
\begin{blocLivre}{cor-24-10}
\begin{corrige}{10}{Regarder par le mauvais bout}
\demarche{Les rôles des lentilles s'échangent : $G'=\dfrac{f'_2}{f'_1}$.}
\begin{questions}
\item La lentille de \qty{20}{mm} devient l'objectif, celle de \qty{900}{mm} l'oculaire.
\item \res{\text{Oui}} : les deux foyers restent confondus.
\item $G'=\dfrac{20}{900}=\res{0{,}022}$ : les objets paraissent \res{\text{45 fois plus petits}}, avec un champ très large.
\item Le \res{\text{judas}} d'une porte, qui montre tout le palier en réduisant l'image.
\end{questions}
\end{corrige}
\end{blocLivre}
\begin{blocLivre}{cor-24-11}
\begin{corrige}{11}{Les jumelles 10 × 50}
\demarche{Largeur du champ $=2\times\qty{1000}{m}\times\tan(\text{demi-champ})$ ; champ apparent $=G\times$ champ réel.}
\begin{questions}
\item Grossissement \res{10} ; diamètre des objectifs \res{\qty{50}{mm}}.
\item $2\times1000\times\tan3{,}25^\circ=\res{\qty{114}{m}}$.
\item $10\times6{,}5^\circ=\res{65^\circ}$.
\item $\dfrac{50}{10}=\res{\qty{5,0}{mm}}$ : inférieur à la pupille la nuit (\qty{7}{mm}), toute la lumière entre dans l'œil ; en
plein jour (\qty{2}{mm}), une partie est perdue, sans gêne puisque la lumière abonde.
\end{questions}
\end{corrige}
\end{blocLivre}
\begin{blocLivre}{cor-24-12}
\begin{corrige}{12}{La grande lunette de Meudon}
\demarche{$G=f'_1/f'_2$ ; $25''=\dfrac{25}{3600}^\circ$ ; pour égaler la Lune, $G=\dfrac{\theta_\text{Lune}}{\theta_\text{Mars}}$.}
\begin{questions}
\item $G=\dfrac{\num{16160}}{25}=\res{646}$.
\item $\theta=\dfrac{25}{3600}\times\dfrac{\pi}{180}=\res{\qty{1,21e-4}{rad}}$ ; $\theta'=646\times\num{1.21e-4}=\qty{0,078}{rad}$, soit \res{4{,}5^\circ}.
\item $G=\dfrac{\num{9.1e-3}}{\num{1.21e-4}}=75$, d'où $f'_2=\dfrac{\num{16160}}{75}=\res{\qty{215}{mm}}$.
\item Un grand objectif est très lourd, se déforme sous son poids et décompose la lumière (aberrations chromatiques) ; un
\res{\text{miroir}}, soutenu par l'arrière, peut être beaucoup plus grand.
\end{questions}
\end{corrige}
\end{blocLivre}
\begin{blocLivre}{cor-24-13}
\begin{corrige}{13}{Les satellites de Jupiter}
\demarche{Angle $=\dfrac{\text{distance Io-Jupiter}}{\text{distance Terre-Jupiter}}$ ; $1'=\dfrac{1}{60}^\circ$.}
\begin{questions}
\item $\theta=\dfrac{\num{4.22e5}}{\num{6.3e8}}=\res{\qty{6,7e-4}{rad}}$, soit \res{2{,}3'}.
\item Le rayon apparent de Jupiter vaut $22{,}5''=0{,}375'$ : écart au bord $\res{1{,}93'}$.
\item $G\geqslant\dfrac{5}{1{,}93}=\res{2{,}6}$ : des jumelles grossissant 10 fois \res{\text{suffisent}} largement.
\item Les satellites \res{\text{tournent autour de Jupiter}} : tous les astres ne tournent pas autour de la Terre, un argument en
faveur du modèle de Copernic.
\end{questions}
\end{corrige}
\end{blocLivre}
\begin{blocLivre}{cor-24-14}
\begin{corrige}{14}{Choisir ses oculaires}
\demarche{$G=f'_1/f'_2$ ; la lentille de Barlow double $f'_1$, donc $G$.}
\begin{questions}
\item Sans Barlow : \res{36} ; \res{75} ; \res{150}. Avec Barlow : \res{72} ; \res{150} ; \res{300}.
\item $2\times80=\res{160}$.
\item 300 est atteint, mais \res{\text{près de deux fois}} au-delà du grossissement utile : l'image serait floue et sombre ; la
publicité est trompeuse.
\item Le grossissement 150, avec l'oculaire de \qty{6}{mm} seul ou celui de \qty{12}{mm} et la Barlow, proche du maximum
utile : Saturne est lumineuse et supporte un fort grossissement.
\end{questions}
\end{corrige}
\end{blocLivre}
\begin{blocLivre}{cor-24-15}
\begin{corrige}{15}{La maquette de lunette}
\demarche{Image intermédiaire dans le plan focal image de l'objectif : $A_1B_1=f'_1\tan\theta$ ; puis $\tan\theta'=A_1B_1/f'_2$.}
\begin{questions}
\item $f'_1=\res{\qty{50}{cm}}$ et $f'_2=\res{\qty{12,5}{cm}}$, distantes de \res{\qty{62,5}{cm}}.
\item L'objet est à l'infini : son image est dans le \res{\text{plan focal image}}, à \qty{50}{cm}, et mesure
$50\times0{,}050=\res{\qty{2,5}{cm}}$.
\item $\theta'=\dfrac{2{,}5}{12{,}5}=\res{\qty{0,20}{rad}}$, d'où $G=\dfrac{0{,}20}{0{,}050}=\res{4}=\dfrac{50}{12{,}5}$.
\item On observe l'image à travers une lunette de visée réglée sur l'infini : elle doit être nette ; ou l'on projette le
faisceau sur un mur lointain, où la tache garde \res{\text{la même taille}} quelle que soit la distance.
\end{questions}
\end{corrige}
\end{blocLivre}
